\chapter{Requerimientos de Negocio}

Este capítulo describe las oportunidades, riesgos y requerimientos de alto nivel que justifican el desarrollo del Sistema de Solicitudes desde la perspectiva del negocio (en este caso, la gestión académica-administrativa de la universidad).

% ============================================================
\section{Oportunidad del Negocio}

El Sistema de Solicitudes representa una oportunidad significativa para la Universidad Autónoma de Zacatecas en los siguientes aspectos:

\begin{enumerate}
    \item \textbf{Mejora de procesos administrativos:} Centralizar y estandarizar las solicitudes reduce los tiempos de atención, elimina la duplicación de esfuerzos y permite al personal enfocarse en resolver trámites en lugar de recopilar información.

    \item \textbf{Transparencia y trazabilidad:} Cada solicitud tendrá un folio único y un historial de seguimiento, lo que permite a los solicitantes conocer el estado de su trámite en todo momento y al personal demostrar que los trámites fueron atendidos.

    \item \textbf{Reducción de errores:} Los formularios dinámicos garantizan que el solicitante envie toda la información y documentación requerida desde el inicio, reduciendo los rechazos por información incompleta.

    \item \textbf{Escalabilidad:} Al ser los tipos de solicitud configurables desde la base de datos (no desde el código), se pueden agregar nuevos tipos de trámite sin necesidad de desarrollo adicional.

    \item \textbf{Generación de métricas:} Por primera vez se podrán obtener datos cuantitativos sobre volumen de solicitudes, tiempos de atención, tipos más demandados y carga de trabajo por area, lo que fácilita la toma de decisiones.

    \item \textbf{Integración con el ecosistema existente:} El sistema se integra como microservicio dentro de la arquitectura existente, aprovechando la autenticación centralizada y los datos del SIGA sin duplicar infraestructura.
\end{enumerate}

% ============================================================
\section{Riesgos del Negocio}

\begin{longtable}{|C{1.2cm}|L{5.5cm}|C{1.5cm}|L{5.5cm}|}
    \hline
    \textbf{ID} & \textbf{Riesgo} & \textbf{Nivel} & \textbf{Mitigación} \\
    \hline
    \endfirsthead
    \hline
    \textbf{ID} & \textbf{Riesgo} & \textbf{Nivel} & \textbf{Mitigación} \\
    \hline
    \endhead
    \hline
    \endfoot

    RI-01 & \textbf{Baja adopción por usuarios.} El personal y los alumnos continuan usando canales informales y no adoptan el sistema. & Alto & Capacitacion a usuarios, interfaz intuitiva, comúnicación institucional sobre el uso obligatorio del sistema. \\
    \hline
    RI-02 & \textbf{Dependencia de servicios externos.} El sistema depende del proveedor de autenticación (JWT) y del SIGA para datos de alumnos. Si alguno falla, el sistema queda parcialmente inoperativo. & Medio & Implementar manejo de errores robusto, cacheo temporal de datos del SIGA, y mensajes claros al usuario cuando un servicio externo no este disponible. \\
    \hline
    RI-03 & \textbf{Resistencia al cambio organizacional.} El personal administrativo puede resistirse a cambiar sus procesos habituales. & Medio & Involucrar al personal desde las etapas de diseño, mostrar beneficios concretos (reducción de carga, trazabilidad). \\
    \hline
    RI-04 & \textbf{Tipos de solicitud no previstos.} Pueden surgir tipos de solicitud con requisitos que el modelo dinámico actual no soporte. & Bajo & El diseño de formularios dinámicos es lo suficientemente flexible para soportar la mayoria de casos. Se puede extender el catálogo de tipos de campo en el futuro. \\
    \hline
    RI-05 & \textbf{Seguridad de datos.} Las solicitudes pueden contener información sensible (documentos personales, comprobantes de pago). & Alto & Comúnicación cifrada (HTTPS), validación de tokens JWT, control de acceso estricto por roles, almacenamiento seguro de archivos. \\
    \hline
    RI-06 & \textbf{Volumen de archivos.} La acumulación de archivos adjuntos puede consumir espacio de almacenamiento significativo con el tiempo. & Bajo & Políticas de retención de archivos, límites de tamaño por archivo, compresión automática. \\
    \hline
\end{longtable}

% ============================================================
\section{Catálogo de Requerimientos de Negocio}

Los requerimientos de negocio representan las necesidades de alto nivel (epicas) que el sistema debe satisfacer.

\begin{longtable}{|C{1.2cm}|L{12.5cm}|}
    \hline
    \textbf{ID} & \textbf{Descripción} \\
    \hline
    \endfirsthead
    \hline
    \textbf{ID} & \textbf{Descripción} \\
    \hline
    \endhead
    \hline
    \endfoot

    RN-01 & Centralizar la gestión de solicitudes académicas y administrativas en una plataforma única accesible via web. \\
    \hline
    RN-02 & Permitir la configuración dinámica de tipos de solicitud, formularios y plantillas sin necesidad de modificar el código fuente. \\
    \hline
    RN-03 & Automatizar el enrutamiento de solicitudes al personal responsable segun el tipo de trámite solicitado. \\
    \hline
    RN-04 & Proporcionar trazabilidad completa del ciclo de vida de cada solicitud con historial de seguimiento. \\
    \hline
    RN-05 & Notificar a las partes involucradas (responsables y solicitantes) sobre eventos relevantes en las solicitudes mediante correo electrónico. \\
    \hline
    RN-06 & Generar documentos oficiales (constancias, etc.) a partir de plantillas configurables con variables dinámicas. \\
    \hline
    RN-07 & Ofrecer métricas y reportes sobre el volumen, tipos y estados de las solicitudes para la toma de decisiones. \\
    \hline
    RN-08 & Integrarse con los servicios existentes de la universidad (autenticación centralizada y SIGA) como un microservicio independiente. \\
    \hline
    RN-09 & Gestiónar excepciones de pago para alumnos mentores mediante una lista administrable. \\
    \hline
\end{longtable}

% ============================================================
\section{Reglas del Negocio}

Las reglas del negocio son restricciónes y políticas del dominio que el sistema debe respetar en todo momento.

\begin{enumerate}[label=\textbf{RG-\arabic*:}, leftmargin=2cm]
    \item Un usuario solo puede tener un rol activo en el sistema (alumno, docente, control escolar, responsable de programa o administrador).

    \item Cada tipo de solicitud debe estar asociado a exactamente un rol responsable de atenderla.

    \item Una solicitud solo puede ser creada por un usuario con rol de alumno o docente.

    \item Los tipos de solicitud disponibles para un alumno pueden diferir de los disponibles para un docente.

    \item Si un tipo de solicitud requiere pago, el solicitante debe adjuntar un comprobante de pago, a menos que el solicitante sea un alumno registrado como mentor activo.

    \item Un alumno mentor queda exento de pago únicamente mientras se encuentre en la lista de mentores activos.

    \item Una solicitud sigue un ciclo de vida estricto: \textit{Creada} $\rightarrow$ \textit{En proceso} $\rightarrow$ \textit{Finalizada} o \textit{Cancelada}. No se puede retroceder de estado.

    \item Una solicitud en estado \textit{Creada} puede ser cancelada por el solicitante. Una vez en estado \textit{En proceso}, solo el personal responsable puede cancelarla o finalizarla.

    \item Cada cambio de estado debe registrarse en el historial de seguimiento con fecha, responsable y observaciones.

    \item El folio de una solicitud es único e inmutable una vez generado.

    \item El personal solo puede ver y atender las solicitudes asignadas a su rol. El administrador puede ver todas las solicitudes.

    \item Las plantillas de documentos solo aplican a ciertos tipos de solicitud y son opcionales.

    \item El sistema no gestióna la autenticación de usuarios; esta se delega a un servicio externo que provee tokens JWT.
\end{enumerate}
